\section{Requirements Specifications}
\subsection{Stakeholder needs and hierarchy}
N01: transport a bag while following one walking golfer. N02: prevent unintended cart motion. N03: identify missing clubs. N04: display distance to green without a phone. N05: remain feasible under budget and schedule. These needs flow into system-level functions (R-F, R-S, R-C, R-G), and then measurable interface/performance requirements (R-IF, R-P).
\begin{table}[ht]\centering\scriptsize\caption{Critical system-level requirements and traceability.}\begin{tabularx}{\linewidth}{p{.13\linewidth}X p{.11\linewidth}p{.13\linewidth}}\toprule ID / priority & Requirement statement & Need & Test\\\midrule
R-F01 / Critical & The cart shall follow one designated golfer at a target separation of 1.5 m with tolerance $\pm0.3$ m in the defined test scenario. & N01 & T03\\
R-S01 / Critical & The cart shall inhibit drive commands within 300 ms of detecting a test obstacle; 200 ms is the design target. & N02 & T04\\
R-S02 / Critical & The cart shall stop or disable motor drive within 100 ms after receipt of a valid emergency-stop command; 50 ms is the design target. & N02 & T05\\
R-S03 / Critical & The cart shall not continue commanded autonomous motion after target-tracking signal loss beyond a validated timeout. & N02 & T06\\
R-C01 / Major & The inventory system shall report present/absent status for 14 designated tagged club slots. & N03 & T07\\
R-G01 / Major & The display shall show the distance from cart GPS position to the selected green with error no greater than 5 m under the defined reference test. & N04 & T08\\
R-P01 / Major & The cart shall limit maximum autonomous cruise speed to 5.0 km/h $\pm0.5$ km/h during acceptance testing. & N01,N02 & T03\\
\bottomrule\end{tabularx}\end{table}
\subsection{Interface, performance and fault requirements}
\begin{table}[ht]\centering\scriptsize\caption{Additional verifiable requirements.}\begin{tabularx}{\linewidth}{p{.16\linewidth}X p{.12\linewidth}}\toprule ID & Requirement statement & Test\\\midrule
R-IF01 & The tracking link shall support a nominal 15 m range ($\pm2$ m) in an unobstructed test area. & T02\\
R-IF02 & The wireless E-stop link shall operate at 20 m minimum range (30 m design target). & T05\\
R-IF03 & The onboard display shall have at least 5 in diagonal size (7 in design target). & T08\\
R-P02 & The cart shall cover 8.5 km or more on one charge under the defined loaded test (10 km design target). & T09\\
R-S04 & The cart shall not energize the drive when E-stop is latched or when controller power is reset. & T05\\
R-C02 & The inventory system shall issue a missing-club warning when a previously detected club is absent during an audit. & T07\\
R-M01 & The prototype shall be constructed within a CAD 600 component budget, counting the planned procurement items. & T10\\
R-X01 / Stretch & The system may log individual shots for future club recommendations. & Not committed\\
\bottomrule\end{tabularx}\end{table}
\subsection{Critical requirement rationale}
R-S01 and R-S02 are safety-critical because a moving loaded cart could contact a golfer or obstacle. Their acceptance measurements refer to detection/command-to-drive-inhibit latency, not stopping distance; braking distance must be recorded separately during T04/T05. R-S03 requires the team to select and document a safe timeout before the test; this parameter is deliberately unresolved. R-G01 requires known reference coordinates and a documented GPS test environment. The 14-slot assumption and range/speed targets originate in the proposal and remain subject to feasibility validation.\par
